% =====================================================================
%  02_arquitectura.tex — Cómo está armado el sistema y por qué
% =====================================================================
\chapter{Arquitectura del sistema}
\label{cap:arquitectura}

\section{Vista general}

\cisa{} es un proceso de Python (contenedor \texttt{cisa-vision}) que arranca la
cámara, lanza dos hilos de visión y un servidor web Flask en el mismo equipo. A
su lado corren dos servicios más: \texttt{cisa-grabacion}, que copia el video
tal cual llega de la cámara, y \texttt{cisa-whatsapp}, que reenvía los avisos a
WhatsApp. La base de datos es un archivo SQLite en el disco de la Jetson.

\begin{figure}[htbp]
  \centering
  \resizebox{\linewidth}{!}{%
  \begin{tikzpicture}[node distance=0.7cm and 0.6cm]

    % Fila 1: adquisición y lectura
    \node[caja] (cam)   at (0,0)    {Cámara\\RTSP};
    \node[caja] (m1)    at (3.8,0)  {M1\\YOLO + ByteTrack};
    \node[caja] (m23)   at (7.6,0)  {M2 + M3\\zona y dígitos};
    \node[caja] (voto)  at (11.2,0) {Votación\\$N$ lecturas iguales};

    % Fila 2: decisión y registro
    \node[caja] (linea) at (11.2,-2.1) {Línea virtual\\entrada / salida};
    \node[caja] (valid) at (7.6,-2.1)  {Validación\\Gemini + Telegram};
    \node[caja] (reg)   at (3.8,-2.1)  {Registro\\SQLite};
    \node[caja] (web)   at (0,-2.1)    {Tablero web\\MJPEG + consultas};

    % Fila 3: servicios paralelos
    \node[cajaGris] (grab)  at (0,-4.2)  {Grabación\\lee el RTSP directo};
    \node[cajaGris] (avisos) at (3.8,-4.2) {Avisos\\Telegram / WhatsApp};

    \draw[flecha] (cam) -- (m1);
    \draw[flecha] (m1) -- node[etiqueta,above]{recortes} (m23);
    \draw[flecha] (m23) -- node[etiqueta,above]{número} (voto);
    \draw[flecha] (voto) -- (linea);
    \draw[flecha] (linea) -- node[etiqueta,above]{cruce} (valid);
    \draw[flecha] (valid) -- (reg);
    \draw[flecha] (reg) -- (web);
    \draw[flecha] (reg) -- (avisos);

    % La grabación lee la misma cámara, sin pasar por la detección
    \draw[flechaGris] (cam.west) -- ++(-2.4,0) |- (grab.west);

    % Recuadro de agrupación
    \begin{scope}[on background layer]
      \node[draw=azulCISA, dashed, rounded corners=4pt,
            fit=(m1)(m23)(voto)(linea)(valid)(reg),
            inner sep=0.3cm,
            label={[font=\scriptsize\bfseries,text=azulCISA]above:%
                   {cisa-vision (un proceso, dos hilos)}}] {};
    \end{scope}
  \end{tikzpicture}%
  }
  \caption{Flujo de un autobús desde la cámara hasta el registro. La grabación
           es independiente de la detección: solo copia el video.}
  \label{fig:flujo}
\end{figure}

\section{Componentes}

\subsection{Cámara (\texttt{logica/camara.py})}

Lee el video por RTSP con decodificación por hardware de la Jetson. Si eso
falla, o si se configura \texttt{CAM\_BACKEND=ffmpeg}, usa FFmpeg. Si
\texttt{RTSP\_URL} apunta a un archivo \texttt{.mp4}, el sistema entra en
\emph{modo prueba}: reproduce el video como si fuera la cámara. La clase
\texttt{CameraStream} mantiene siempre el cuadro más reciente en
\texttt{logica/estado.py}, de donde lo toman la detección, el video en vivo y
las capturas.

\subsection{Detección M1 (\texttt{ml/hilos.py}, \texttt{HiloM1})}

YOLO detecta autobuses y ByteTrack les asigna un identificador que se conserva
mientras el camión está en cuadro. El hilo está limitado por
\texttt{YOLO\_FPS\_MAX} (0 = sin límite, procesa cada cuadro nuevo). Su salida
son recortes de cada autobús que se pasan al segundo hilo por una cola.

\subsection{Lectura M2 + M3 (\texttt{ml/hilos.py}, \texttt{HiloM2M3})}

\begin{itemize}
  \item \textbf{M2} localiza la zona del número y la clasifica como
        \emph{delantera}, \emph{lateral} o \emph{trasera}. Cada zona detectada
        se asigna al autobús que la contiene.
  \item \textbf{M3} lee los dígitos de esa zona.
\end{itemize}

Un autobús necesita \textbf{2 lecturas iguales} (\texttt{N\_VOTOS}) para que su
número quede confirmado. El hilo respeta \texttt{M2\_FPS\_MAX} y dos topes por
cuadro: \texttt{MAX\_ZONAS\_POR\_BUS} (2 zonas por camión) y
\texttt{MAX\_LECTURAS\_M3} (4 lecturas en total). Gracias a esos topes, la carga
de la GPU no crece aunque haya muchos camiones a la vez.

\subsection{Línea virtual (\texttt{ml/linea\_virtual.py})}

La línea es un segmento definido por dos puntos en la imagen
(\texttt{LINEA\_P1}, \texttt{LINEA\_P2}). El sistema calcula de qué lado está el
centro del autobús y lo compara con \texttt{LADO\_ENTRADA}. Alrededor de la
línea hay un margen (\texttt{MARGEN\_LINEA\_PX}) en el que el lado
\emph{no cambia}: así el temblor de la detección con el camión detenido sobre la
línea no se cuenta como cruce.

\subsection{Registro de un solo evento por cruce}

Cuando un número confirmado cruza, el sistema no registra de inmediato. Espera
a que el camión termine de pasar (\texttt{ESPERA\_FIN\_SEG} sin verlo, con un
máximo de \texttt{ESPERA\_MAX\_SEG}) y entonces elige el mejor número leído.
Las reglas son:

\begin{enumerate}
  \item Si dos números compiten y uno es otro con dígitos de menos
        (\texttt{504} dentro de \texttt{5014}), gana el más largo.
  \item Si el número más leído \textbf{no está en el padrón} y se leyó al menos
        \texttt{VOTOS\_FUERA\_PADRON} veces igual, se registra tal cual como
        \emph{fuera del padrón}. No se ``corrige'' a un número parecido
        (\texttt{2508} no se convierte en \texttt{5008}).
  \item Un mismo número no se vuelve a registrar antes de
        \texttt{DEBOUNCE\_SEG} (60 s por defecto).
  \item Si la unidad regresa antes de \texttt{MANIOBRA\_SEG}, el movimiento se
        considera maniobra y se descarta.
\end{enumerate}

\subsection{Validación (\texttt{avisos/gemini.py}, \texttt{avisos/telegram.py})}

Antes de dar por bueno un cruce, Gemini revisa la foto. Si hay duda, se manda
la imagen a los validadores por Telegram para que una persona confirme o
corrija. En pruebas se apaga con \texttt{SIN\_TELEGRAM=1}. La validación con
Gemini se controla aparte con \texttt{VALIDAR\_CON\_GEMINI}.

\subsection{Persistencia (\texttt{metricas/})}

La base SQLite guarda eventos y estancias; el padrón \textbf{no} se guarda como
fuente, se sincroniza desde \path{padron.csv} en cada arranque. El detalle del
esquema está en el capítulo~\ref{cap:datos}.

\subsection{Web (\texttt{web/})}

Flask sirve el tablero, el video en vivo por MJPEG, las fotos, las grabaciones y
una API JSON. Las rutas solo \emph{leen}: no contienen lógica de visión ni de
almacenamiento. El detalle de los endpoints está en el
capítulo~\ref{cap:referencia}.

\subsection{Grabación (\texttt{scripts/grabar.sh})}

Copia el stream de la cámara sin recodificar (Full HD H.265), un archivo por
hora. Conserva 72 horas y rota los archivos viejos. Corre como servicio aparte
para que una falla en la detección no detenga la grabación, y viceversa.

\subsection{WhatsApp (\texttt{whatsapp/})}

Servicio Node.js con Baileys que mantiene una sesión de WhatsApp Web y reenvía
los mismos avisos del grupo de Telegram a \texttt{WHATSAPP\_NUMERO}. No usa la
API oficial. La sesión vive en \path{whatsapp/sesion/} y equivale a la
contraseña de la cuenta: no se comparte ni se sube a Git.

\section{Mapa de módulos}

\begin{longtable}{@{}L{4.3cm}L{4.3cm}L{5.7cm}@{}}
  \toprule
  \textbf{Archivo} & \textbf{Qué es} & \textbf{Responsabilidad} \\
  \midrule
  \endhead
  \path{config.py} & Configuración & Único archivo que lee el \texttt{.env}; valida
    variables, modelos, carpetas y padrón al arrancar. \\
  \path{pipe_stream_v2.py} & Arranque & Crea la cámara, lanza los hilos y el
    servidor web. \\
  \path{logica/camara.py} & Adquisición & Lectura RTSP, respaldo FFmpeg, modo
    prueba con \texttt{.mp4}. \\
  \path{logica/estado.py} & Estado compartido & Último cuadro, estado del
    pipeline y referencias de runtime. \\
  \path{ml/modelos.py} & Modelos & Carga y calentamiento de los motores
    TensorRT; helpers de zona. \\
  \path{ml/hilos.py} & Hilos de visión & \texttt{HiloM1} (detección) y
    \texttt{HiloM2M3} (lectura). \\
  \path{ml/preprocesamiento.py} & Preproceso & Ajustes a la imagen antes de leer
    los dígitos. \\
  \path{ml/reconocimiento.py} & Reconocimiento & Recortes con margen, lectura de
    dígitos, validación del número, guardado de crops. \\
  \path{ml/linea_virtual.py} & Cruce & Lado de la línea, distancia, evaluación
    del cruce y reasociación de números. \\
  \path{metricas/schema.sql} & Esquema & Tablas, índices y vistas de la base. \\
  \path{metricas/storage.py} & Persistencia & Conexión por hilo, migración,
    sincronización del padrón, altas y consultas. \\
  \path{metricas/metricas.py} & Agregados & Cálculos que alimentan el tablero. \\
  \path{avisos/telegram.py} & Avisos & Envío de fotos y texto a Telegram y a los
    validadores. \\
  \path{avisos/gemini.py} & Validación & Revisión automática de la foto. \\
  \path{avisos/whatsapp.py} & Avisos & Puente hacia el servicio de WhatsApp. \\
  \path{web/servidor.py} & Flask & Instancia de la aplicación. \\
  \path{web/rutas.py} & HTTP & Todas las rutas: API, páginas, MJPEG, archivos. \\
  \path{templates/}, \path{static/} & Interfaz & Páginas HTML, CSS y JavaScript. \\
  \path{whatsapp/puente.js} & WhatsApp & Sesión de WhatsApp Web (Node + Baileys). \\
  \path{scripts/} & Operación & Cambio de fuente, grabación, monitoreo y red. \\
  \bottomrule
  \caption{Módulos del repositorio y su responsabilidad.}
  \label{tab:modulos}
\end{longtable}

\section{Concurrencia y rendimiento}

El sistema está pensado para una sola Jetson compartida por todo el pipeline:

\begin{itemize}
  \item \textbf{Un cuadro, muchos consumidores.} La cámara publica el cuadro más
        reciente; nadie guarda copias de más.
  \item \textbf{Dos hilos con límites independientes.} M1 y M2/M3 tienen su
        propio tope de cuadros por segundo (\texttt{YOLO\_FPS\_MAX},
        \texttt{M2\_FPS\_MAX}). Bajar el segundo alivia la GPU sin dejar de
        detectar autobuses.
  \item \textbf{Topes por cuadro.} M3 nunca hace más de
        \texttt{MAX\_LECTURAS\_M3} lecturas por cuadro, así que el consumo no
        depende del número de camiones en imagen.
  \item \textbf{Relectura selectiva.} Un camión con número ya confirmado que no
        ha cruzado se vuelve a leer cada \texttt{RELECTURA\_SEG} segundos, para
        no gastar GPU en lo mismo.
  \item \textbf{Sin base de datos en el camino crítico de la imagen.} La
        escritura ocurre una sola vez por cruce, no por cuadro.
\end{itemize}

\begin{aviso}[title={Antes de subir los FPS}]
Subir \texttt{M2\_FPS\_MAX} mejora la lectura cuando los camiones pasan rápido,
pero calienta la Jetson. Mide el consumo y la temperatura antes de dejarlo
permanente. Las cifras medidas están en el capítulo~\ref{cap:operacion}.
\end{aviso}

\section{Decisiones de diseño}

\begin{description}[style=nextline]
  \item[Un solo \texttt{config.py}.] Ningún otro archivo usa
        \texttt{os.getenv}. Así se sabe de un vistazo qué variables existen y
        qué pasa si falta una.
  \item[El padrón es un archivo, no una tabla.] Se edita \path{padron.csv} y se
        reinicia. Evita dos fuentes de verdad para las unidades válidas.
  \item[Votación en vez de una lectura.] Una sola lectura confunde \texttt{5}
        con \texttt{6} o \texttt{8} con \texttt{3}; dos lecturas iguales bajan
        mucho el error sin frenar el paso.
  \item[Margen de línea.] El detector tiembla un par de píxeles con el camión
        detenido. Sin margen, un camión esperando en la puerta generaría
        entradas y salidas en cadena.
  \item[Esperar a que termine de pasar.] Registrar al instante obliga a decidir
        con la primera lectura. Esperar unos segundos permite elegir el mejor
        número y evita dos registros del mismo cruce.
  \item[Fuera del padrón es un estado válido.] Es mejor registrar ``no
        reconocido'' con evidencia que inventar una unidad parecida.
\end{description}
